

\chapter{Resultados e Discussões}
\label{cha:resultados-e-discussoes}

Para investigar o desempenho do \ac{aog}, foram feitos testes em inversões com dados sintéticos e dados reais. Utilizou-se a ferramenta MATLAB nos experimentos.

Obteve-se os dados das sondagens elétricas verticais de Schlumberger dos modelos sintéticos baseados em \citeonline{Sen1993} sendo aplicados 5\% de ruído gaussiano pseudoaleatório em ambos os dados sintéticos, para simular o ruído dos dados observados ao realizar medidas em campo. 

Foram feitas 10 inversões para cada um dos modelos sintéticos. Cada inversão foi realizada pelo \ac{aog} com 1000 iterações e uma população de 100 gafanhotos para o modelo com três camadas e 100 gafanhotos para o modelo de 5 camadas. Com o resultado das inversões, foi obtida a solução média e desvio padrão, gerando diferentes modelos médios.

Já no modelo real, baseia-se em \citeonline{bhattacharya1968direct}, foram feitas 10 inversões com 1000 iterações e uma população de 150 gafanhotos para o modelo real de três camadas. Com o resultado das inversões, foi obtida a solução média e desvio padrão, gerando diferentes modelos médios.

Também serão exibidas figuras que mostram as curvas de convergência da função \textit{fitnnes} das inversões. É importante ressaltar que cada figura mostra um exemplo de uma única inversão dentre as iterações.


% Dados sintéticos foram gerados para o modelo usando o código de modelagem direta e pontos dados por círculos mostrados na figura (ref). Para o \ac{aog}, usamos limites de pesquisa ub e lb na Tabela(ref).
\section{Inversão de dados sintéticos em três camadas (Modelo 1)}
\label{sec:inversao_sinteticos}

O desempenho do \ac{aog} na inversão em uma \ac{sev} é primeiramente investigado em dados sintéticos gerados a partir de um modelo de três camadas. Na tabela \ref{tab:tab_sint_3} contém os valores verdadeiros e espaço de busca utilizado na inversão. $\rho_i$ e $E_i$ correspondem, respectivamente, aos valores de resistividade e espessura da camada $i$. A última camada tem espessura infinita. O valor mínimo e o valor máximo correspondem aos limites do espaço de busca.

\begin{table}[H]
    \centering
    \caption{Parâmetros para modelo de três camadas.}
    \begin{tabular}{@{}ccccc@{}}
        \toprule
        \textbf{Parâmetros}  & \textbf{Valores Verdadeiros} & \textbf{Mínimo}      & \textbf{Máximo}      \\ \midrule
        $\rho_1(\Omega$m)    & 10                           & 5                    & 15                   \\
        $\rho_2(\Omega$m)    & 390                          & 15                   & 500                  \\
        $\rho_3(\Omega$m)    & 10                           & 1                    & 20                   \\
        $E_1($m)             & 10                           & 1                    & 20                   \\
        $E_2($m)             & 250                          & 100                  & 500                  \\
        \multicolumn{1}{l}{} & \multicolumn{1}{l}{}         & \multicolumn{1}{l}{} & \multicolumn{1}{l}{} \\ \bottomrule
    \end{tabular}
    \label{tab:tab_sint_3}
\end{table}


A solução média e desvio padrão das 10 inversões entre o modelo verdadeiro e o modelo estimado em três camadas com ruído estão descritos na tabela \ref{tab:tab_media_3}.



\begin{table}[H]
    \center
   \caption{Média e desvio padrão de 10 inversões.}
    \begin{tabular}{@{}ccc@{}}
        \toprule
        \textbf{Parâmetros}  & \textbf{Valor Médio} & \textbf{Desvio Padrão} \\ \midrule
        $\rho_1(\Omega$m)    & 10.07                & 1.10                   \\
        $\rho_2(\Omega$m)    & 371.41               & 12.12                  \\
        $\rho_3(\Omega$m)    & 9.10                 & 6.0                    \\
        $E_1($m)             & 10.01                & 1.13                   \\
        $E_2($m)             & 263.39               & 18.18                  \\
        \multicolumn{1}{l}{} & \multicolumn{1}{l}{} & \multicolumn{1}{l}{}   \\
        \bottomrule
    \end{tabular} 
    
    \label{tab:tab_media_3}
\end{table}

A figura \ref{fig:mod_1_figure} mostra a curva de resistividade média com 100 ganhotos e 1000 interações. Nota-se que o \ac{aog} recuperou muito entre dados da resistividade aparente observada (círculos vermelhos) e resistividade aparente calculada teoricamente (linha sólida azul).


\begin{figure}[H]
    \centering

    \includegraphics[width=1\linewidth]{lib/fig/result/fig1_mod1.eps}
    \caption{ Resultados da inversão do modelo sintético de 3 camadas usando AOG, comparação entre dados da resistividade aparente observada (círculos vermelhos) e resistividade aparente calculada teoricamente (linha sólida azul) para o melhor ajuste do modelo AOG.}
    \label{fig:mod_1_figure}
\end{figure}

Na figura \ref{fig:mod_1_modelo_figure} mostra o modelo de camadas obtidos com a inversão utilizando o \ac{aog}. A Camada um é bem resolvida. A segunda e terceira camada também apresenta boa resolução, embora o modelo invertido superestime a resistividade dessa camada.

\begin{figure}[H]
    \centering

    \includegraphics[width=1\linewidth]{lib/fig/result/fig2_mod1.eps}
    \caption{Modelo verdadeiro (linha sólida azul) e modelo médio estimado para o AOG (linha tracejada vermelha).}

    \label{fig:mod_1_modelo_figure}
\end{figure}

A Figura \ref{fig:mod_1_modelo_convergence} representa uma curva típica de uma inversão, percebe-se que a curva converge bem para a solução da inversão.
\begin{figure}[H]
    \centering
    \includegraphics[width=1\linewidth]{lib/fig/result/fig3_mod1.eps}
    \caption{Curva típica de convergência do modelo sintético de 3 camadas usando o AOG de uma inversão.}
    \label{fig:mod_1_modelo_convergence}
\end{figure}

\section{Inversão de dados sintéticos em cinco camadas (Modelo 2)}
\label{sec:inversao_5_sintetico}
O desempenho do \ac{aog} é também investigado em uma \ac{sev} de cinco camadas. Os espaços de busca (valores máximo e mínimo) dos parâmetros do modelo são mostrados na tabela \ref{tab:tab_sint_5}.


\begin{table}[H]
    \centering
    \caption{Parâmetros para modelo de cinco camadas.}
    \begin{tabular}{@{}ccccc@{}}
        \toprule
        \textbf{Parâmetros}  & \textbf{Valores Verdadeiros} & \textbf{Mínimo}      & \textbf{Máximo}      \\ \midrule
        $\rho_1(\Omega$m)    & 10                           & 1                    & 20                   \\
        $\rho_2(\Omega$m)    & 50                           & 30                   & 60                   \\
        $\rho_3(\Omega$m)    & 100                          & 50                   & 150                  \\
        $\rho_4(\Omega$m)    & 20                           & 10 0                 & 40                   \\
        $\rho_5(\Omega$m)    & 400                          & 200                  & 600                  \\
        $E_1($m)             & 2                            & 0.5                  & 5                    \\
        $E_2($m)             & 15                           & 10                   & 20                   \\
        $E_3($m)             & 20                           & 10                   & 30                   \\
        $E_4($m)             & 25                           & 10                   & 50                   \\
        \multicolumn{1}{l}{} & \multicolumn{1}{l}{}         & \multicolumn{1}{l}{} & \multicolumn{1}{l}{} \\ \bottomrule
    \end{tabular}
    \label{tab:tab_sint_5}
\end{table}

 A solução média e desvio padrão das 10 inversões entre o modelo verdadeiro e o modelos estimado da inversão em cinco camadas com ruído estão descritos na Tabela \ref{tab:tab_media_5}.

\begin{table}[H]
    \center
    \caption{Média e desvio padrão de 10 inversões.}
    \begin{tabular}{@{}ccc@{}}
        \toprule
        \textbf{Parâmetros}  & \textbf{Valor Médio} & \textbf{Desvio Padrão} \\ \midrule
        $\rho_1(\Omega$m)    & 11.93                & 4.10                   \\
        $\rho_2(\Omega$m)    & 55.76                & 7.62                   \\
        $\rho_3(\Omega$m)    & 12.11                & 32.39                  \\
        $\rho_4(\Omega$m)    & 19.10                & 6.01                   \\
        $\rho_5(\Omega$m)    & 362.05               & 4.38                   \\
        $E_1($m)             & 3.08                 & 1.44                   \\
        $E_2($m)             & 14.74                & 4.00                   \\
        $E_3($m)             & 16.81                & 7.57                   \\
        $E_4($m)             & 23.58                & 8.58                   \\
        \multicolumn{1}{l}{} & \multicolumn{1}{l}{} & \multicolumn{1}{l}{}   \\
        \bottomrule
    \end{tabular}
    \label{tab:tab_media_5}
\end{table}


 Com 100 gafanhotos em 1000 iterações nota-se que a \ac{aog} recuperou bem a parâmetros da inversão do modelo verdadeiro (linha azul sólida)  em comparação como modelo estimado (linhas tracejada vermelha) (Figura \ref{fig:mod_2_modelo_figure}).
 
\begin{figure}[H]
    \centering

    \includegraphics[width=1\linewidth]{lib/fig/result/fig1_mod2.eps}
    \caption{Resultados da inversão do modelo sintético de 5 camadas usando AOG, comparação entre dados da resistividade aparente observada (círculos vermelhos) e resistividade aparente calculada teoricamente (linha sólida azul) para o melhor ajuste do modelo AOG.}
    % \source{O Autor}
    \label{fig:mod_2_figure}
\end{figure}


Na figura \ref{fig:mod_2_modelo_figure} mostra o modelo de camadas obtidos com a inversão utilizando o \ac{aog}. A três camadas apresentam uma  resolução mediana, embora o modelo invertido superestime a resistividade dessa camada, é aceitável.

\begin{figure}[H]
    \centering

    \includegraphics[width=1\linewidth]{lib/fig/result/fig2_mod2.eps}
    \caption{Modelo de camadas}

    \label{fig:mod_2_modelo_figure}
\end{figure}

A Figura \ref{fig:mod_2_modelo_convergence} representa uma curva típica de uma inversão, percebe-se que a curva converge bem para a solução da inversão.
\begin{figure}[H]
    \centering

    \includegraphics[width=1\linewidth]{lib/fig/result/fig3_mod2.eps}
    \caption{Curva típica de convergência do modelo sintético de 5 camadas usando o AOG de uma inversão.}
    \label{fig:mod_2_modelo_convergence}
\end{figure}

\section{Inversão de dados reais em três camadas (Modelo 3)}
\label{sec:inversao_3_real}


O  \ac{aog}, na inversão em dados reais, foi também investigado nesse trabalho. A Figura \ref{fig:mod_3_figure} ilustra o resultado de inversão do \ac{aog} em uma curva de resistividade do modelo 3 em comparação com o resultado do \ac{aog}. Os espaços de busca (valores máximo e mínimo) dos parâmetros do modelo são mostrados na Tabela \ref{tab:tab_real_3}.

\begin{table}[H]
    \centering
    \caption{Parâmetros para modelo real de três camadas.}
    \begin{tabular}{@{}ccc@{}}
        \toprule
        \textbf{Parâmetros}   & \textbf{Mínimo}      & \textbf{Máximo}      \\ \midrule
        $\rho_1(\Omega$m)                             & 10                    & 100                   \\
        $\rho_2(\Omega$m)                               & 1                  & 10                  \\
        $\rho_3(\Omega$m)                              & 5                    & 20                   \\
        $E_1($m)                                         & 1                    & 5                   \\
        $E_2($m)                                        & 10                  & 25                 \\
        \multicolumn{1}{l}{}         & \multicolumn{1}{l}{} & \multicolumn{1}{l}{} \\ \bottomrule
    \end{tabular}
    \label{tab:tab_real_3}
\end{table}




\begin{figure}[H]
    \centering

    \includegraphics[width=1\linewidth]{lib/fig/result/fig1_mod3.eps}
    \caption{ Resultados da inversão do modelo real de 3 camadas usando o AOG, comparação entre dados da resistividade aparente observada (círculos vermelhos) e resistividade aparente calculada teoricamente (linha sólida azul) para o melhor ajuste do modelo AOG.}
    \label{fig:mod_3_figure}
\end{figure}


Nota-se que a \ac{aog} recuperou muito bem a parâmetros da inversão do modelo verdadeiro (linha azul sólida)  em comparação como modelo estimado (linhas tracejada vermelha) (Figura \ref{fig:mod_3_modelo_figure}). A solução média das 10 inversões entre o modelo verdadeiro e o modelos estimado da inversão em três camadas estão descritos na Tabela \ref{tab:tab_media_3_real}. Com 150 gafanhotos em 1000 iterações o \ac{aog} se ajusta aos dados de resistividade aparente.



\begin{table}[H]
    \center
   \caption{Média e desvio padrão de 10 inversões.}
    \begin{tabular}{@{}ccc@{}}
        \toprule
        \textbf{Parâmetros}  & \textbf{Valor Médio} & \textbf{Desvio Padrão} \\ \midrule
        $\rho_1(\Omega$m)    & 32.82                & 0.82                   \\
        $\rho_2(\Omega$m)    & 2.67               & 1.01                  \\
        $\rho_3(\Omega$m)    & 14.00                 & 1.23                   \\
        $E_1($m)             & 3.83                & 0.30                   \\
        $E_2($m)             & 19.74               & 5.48                  \\
        \multicolumn{1}{l}{} & \multicolumn{1}{l}{} & \multicolumn{1}{l}{}   \\
        \bottomrule
    \end{tabular} 
  
    \label{tab:tab_media_3_real}
\end{table}


\begin{figure}[H]
    \centering

    \includegraphics[width=1\linewidth]{lib/fig/result/fig2_mod3.eps}
    \caption{Modelo verdadeiro (linha sólida azul) e modelo médio estimado para o AOG (linha tracejada vermelha).}

    \label{fig:mod_3_modelo_figure}
\end{figure}

A Figura \ref{fig:mod_3_modelo_convergence} representa uma curva típica de uma inversão. A curva convergiu bem para a solução da inversão.
\begin{figure}[H]
    \centering
    \includegraphics[width=1\linewidth]{lib/fig/result/fig3_mod3.eps}
    \caption{Curva típica de convergência do modelo real de 3 camadas usando o AOG de uma inversão.}
    \label{fig:mod_3_modelo_convergence}
\end{figure}


